%% ============================================================================
%% Fichier  : annexes/Ann_D_textes.tex
%% Titre    : Textes Juridiques et Normatifs
%% Version  : 0.1 -- novembre 2026
%% ============================================================================
%% Textes lus dans le dossier Normes : décret n° 2013/287
%% (texte_organique_consupe.pdf) ; loi n° 2010/012. Les autres textes sont
%% décrits sans citation.
%% ============================================================================

\chapter{Textes Juridiques et Normatifs}
\label{ann:textes}

Les extraits ci-dessous sont recopiés de la version des textes disponible
dans le dossier des normes du cours. L'auditeur les vérifie sur la version
officielle avant de les citer dans un rapport.

\section*{D.1 Décret n° 2013/287 du 4 septembre 2013 portant organisation
des services du Contrôle supérieur de l'État}

\begin{boxjuris}
\textbf{Article~1er}~: \textit{« (1) Les services du Contrôle supérieur de
l'Etat relèvent de l'autorité directe du Président de la République, dont
ils reçoivent les instructions et à qui ils rendent compte. (2) Ils sont
dirigés par un ministre délégué à la présidence de la République. »}

\medskip
\textbf{Article~2, alinéa~1} (début)~: \textit{« Les services du Contrôle
supérieur de l'Etat constituent l'institution supérieure de contrôle des
finances publiques (ISC) du Cameroun. Ils sont chargés de l'audit
externe. »}

\medskip
\textbf{Article~2, alinéa~3}~: \textit{« Dans l'exercice de leurs
attributions, les services du Contrôle supérieur de l'Etat effectuent~:
un contrôle de conformité et de régularité~; un contrôle financier~; un
contrôle de performance~; un contrôle de l'environnement~; un contrôle des
systèmes d'information. »}

\medskip
\textbf{Article~2, alinéa~4}~: \textit{« Les contrôles sont effectués par
les vérificateurs du Contrôle supérieur de l'Etat sur pièces et sur place.
A cet effet, et nonobstant les contrôles effectués par les autres organes
de contrôle, les responsables des structures auditées sont tenus de
présenter les pièces relatives à la période auditée. »}

\medskip
\textbf{Article~3, alinéa~1}~: \textit{« Les services du Contrôle supérieur
de l'Etat concourent à la sanction des ordonnateurs et des gestionnaires
des deniers publics, dans les conditions prévues par les lois et
règlements en vigueur. A ce titre, le ministre délégué en charge du
Contrôle supérieur de l'Etat préside le conseil de discipline budgétaire
et financière, dont l'organisation et le fonctionnement sont fixés par un
texte particulier. »}

\medskip
\textbf{Article~5}~: \textit{« (1) Les structures d'audit interne
transmettent aux services du Contrôle supérieur de l'Etat, à la diligence
des chefs de départements ministériels concernés, un exemplaire de chacun
des rapports rédigés à la suite de toutes les interventions menées au
cours de l'exercice, ainsi qu'un exemplaire du rapport annuel d'activités.
(2) Après examen des rapports visés à l'alinéa 1er ci-dessus, les services
du Contrôle supérieur de l'Etat peuvent, le cas échéant, saisir le
conseil de discipline budgétaire et financière. »}

\medskip
\textbf{Article~27, alinéa~2} (extrait)~: la Division des inspections et
de contrôle des administrations publiques comprend notamment \textit{« la
section des administrations financières, économiques, sociales et des
systèmes d'informations »}.

\medskip
\textbf{Article~28, alinéa~1}~: les sections sont chargées, chacune dans
son domaine de compétence, \textit{« de la préparation et de l'exécution
des opérations de vérification~; du suivi et de l'évaluation des
activités des structures de contrôle interne~; de la tenue de la
documentation technique et du fichier. »}
\end{boxjuris}

\section*{D.2 Loi n° 2010/012 du 21 décembre 2010 relative à la
cybersécurité et à la cybercriminalité}

\begin{boxjuris}
\textbf{Article~4, point~9}~: \textit{« Audit de sécurité~: examen
méthodique des composantes et des acteurs de la sécurité, de la
politique, des mesures, des solutions, des procédures et des moyens mis
en oeuvre par une organisation, pour sécuriser son environnement,
effectuer des contrôles de conformité, des contrôles d'évaluation de
l'adéquation des moyens (organisationnels, techniques, humains,
financiers) investis au regard des risques encourus, d'optimisation, de
rationalité et de performance. »}

\medskip
\textbf{Article~7, alinéa~2} (extrait)~: l'Agence est notamment chargée
\textit{« de s'assurer de la régularité, de l'effectivité des audits de
sécurité des systèmes d'information suivant les normes en la matière, des
organismes publics et des autorités de certification »}.

\medskip
\textbf{Article~13, alinéa~1}~: \textit{« Sont soumis à un audit de
sécurité obligatoire, les réseaux de communications électroniques et les
systèmes d'information, des opérateurs, des autorités de certification et
des fournisseurs de services de communications électroniques. »}

\medskip
\textbf{Article~32}~: \textit{« (1) Les réseaux de communications
électroniques et les systèmes d'information sont soumis à un régime
d'audit de sécurité obligatoire et périodique de leurs systèmes de
sécurité par l'Agence. (2) L'audit de sécurité et les mesures d'impact de
gravité sont effectués chaque année ou lorsque les circonstances
l'exigent. (3) Les rapports d'audit sont confidentiels et adressés au
Ministre chargé des Télécommunications. »}
\end{boxjuris}

\section*{D.3 Textes décrits sans citation}

\begin{tableaucours}{L{4.8cm}L{7.4cm}}{%
    \tableauHeader{Texte} &
    \tableauHeader{Usage dans le cours}
}
Loi n°~2018/012 du 11~juillet~2018 portant régime financier de l'État et
des autres entités publiques & Exécution de la dépense, séparation
ordonnateur / comptable (Ch.~\ref{chap:sci}) \\
Texte organisant la Chambre des Comptes de la Cour suprême & Contrôle des
comptes~; les SI par extension (Ch.~\ref{chap:institutions}) \\
Loi n°~2024/017 du 23~décembre~2024 relative à la protection des données
personnelles & Données des agents, contribuables, usagers
(Ch.~\ref{chap:gouvernance}) \\
\end{tableaucours}

\section*{D.4 Normes}

\begin{tableaucours}{L{4.4cm}L{7.8cm}}{%
    \tableauHeader{Norme} &
    \tableauHeader{Passages utilisés}
}
COSO~2013 (traduction IFACI) & Définition, composantes et principes,
efficacité, déficiences (§§~138 à~140) \\
INTOSAI GUID~5100 & Audit des SI autonome ou intégré~; treize domaines \\
Manuel d'audit de conformité de l'IDI & Relation tripartite, critères,
documentation \\
Normes de l'IIA (2017, 2024) & Exigences de l'audit interne \\
COBIT~2019, ISO/IEC~27001 et 27002 & Sources de critères (résumés, aucune
reproduction) \\
\end{tableaucours}
